\chapter{Configuration and task spaces}
\label{chap:lecture11}
Direct kinematics describes the result of a robot configuration. Planning
asks how to move from a starting configuration to a configuration that
accomplishes a task. A mobile robot may need to reach a doorway, while a
manipulator may need to place its tool at an object. Both problems require
a description of the robot's possible configurations and of the desired
outcome. These descriptions need not use the same coordinates.

\section{A configuration describes the robot's geometry}
A \emph{configuration} specifies the placement of the robot's bodies,
subject to its mechanical constraints. A configuration vector $q$ collects
independent coordinates sufficient to determine that placement, given the
robot's geometry and base-frame definition.

For a rigid mobile robot moving in a plane, one choice is
\[
 q=\begin{bmatrix}x_R^W&y_R^W&\theta_R^W\end{bmatrix}\trans.
\]
The position and heading determine its footprint in the environment.
For a fixed-base planar RR arm, the corresponding description is
\[
 q=\begin{bmatrix}q_1&q_2\end{bmatrix}\trans.
\]
Here the coordinates are joint angles rather than world positions.
Nevertheless, knowing them determines the placement of both links and
the tool through forward kinematics.

\begin{figure}[htbp]
\centering\begin{tikzpicture}[scale=.85]
\draw[axis] (0,0)--(4.2,0) node[right] {$x_W$};
\draw[axis] (0,0)--(0,3.2) node[above] {$y_W$};
\draw[projection] (2.4,0) node[below] {$x_R^W$}--(2.4,1.7)--(0,1.7) node[left] {$y_R^W$};
\begin{scope}[shift={(2.4,1.7)},rotate=30]
\draw[fill=gray!15,thick] (-.7,-.45) rectangle (.7,.45);
\draw[line width=3pt] (-.3,-.5)--(.3,-.5) (-.3,.5)--(.3,.5);
\draw[axis,red!75!black] (0,0)--(1.3,0) node[right] {$x_R$};
\end{scope}
\draw[dashed] (2.4,1.7)--(3.7,1.7);
\draw[->] (3.4,1.7) arc (0:30:1);
\node[right] at (3.65,1.65) {$\theta_R^W$};
\node at (2,-.9) {$q=(x_R^W,y_R^W,\theta_R^W)\trans$};
\begin{scope}[shift={(7,0)}]
\coordinate (O) at (0,0); \coordinate (E) at (30:2.2);
\coordinate (P) at ($(E)+(100:1.8)$);
\draw[axis] (O)--(3.2,0) node[right] {$x_0$};
\draw[axis] (O)--(0,3.2) node[above] {$y_0$};
\draw[line width=8pt,gray!45] (O)--(E)--(P);
\draw[thick] (O)--(E)--(P);
\foreach \p in {O,E,P} \filldraw[fill=white] (\p) circle (.09);
\draw[->] (.8,0) arc (0:30:.8) node[right] {$q_1$};
\draw[dashed] (E)--++(30:1.2);
\draw[->] ($(E)+(30:.7)$) arc (30:100:.7) node[above] {$q_2$};
\node[above] at (P) {tool};
\node at (1.5,-.9) {$q=(q_1,q_2)\trans$};
\end{scope}
\end{tikzpicture}

\caption{Two configuration descriptions: world position and heading for
a planar mobile robot, and relative joint angles for a fixed-base RR arm.
Each point in the corresponding configuration space represents an entire
robot placement, not just the location of its tool or reference point.}
\label{fig:l11-configurations}
\end{figure}

The \emph{configuration space}, denoted $\mathcal C$, is the set of
configurations admitted by the kinematic model and any specified joint
limits. Environmental collision constraints will be treated separately
through a free subset of $\mathcal C$.

Configuration is not generally the complete dynamic state. Velocities,
and sometimes actuator or other internal states, are also needed to
predict future motion from applied inputs. Here the question is geometric:
where are the robot's parts? The dimension of $\mathcal C$ is the number
of independent configuration degrees of freedom. A coordinate vector is
a representation of this space; angular coordinates require care at their
chosen wrap boundaries.

\section{The task selects the output coordinates}
A \emph{task space} describes the quantities relevant to the requested
outcome. For example:
\begin{center}
\begin{tabular}{p{.44\linewidth}p{.46\linewidth}}
\toprule
Task & Possible task coordinates\\
\midrule
Reach a location in a plane & Position $(x,y)$\\
Hold a point in three-dimensional space & Position $(x,y,z)$\\
Maintain altitude & Height $z$\\
Keep a container upright & Tilt, represented locally by roll and pitch\\
Place and orient a tool & Full position and orientation\\
\bottomrule
\end{tabular}
\end{center}
A position-only task does not prescribe final orientation. Keeping a
container upright need not prescribe yaw, although carrying liquid may
also require limits on acceleration. A more detailed task may constrain
several points or axes on the robot. Collision avoidance must account for
the robot's full geometry even if the task output contains only tool
position.

Write the task output as $\xi$ and the task space as $\mathcal X$.
The robot model defines a task map
\begin{equation}
 \xi=f(q),\qquad f:\mathcal C\longrightarrow\mathcal X.
 \label{eq:l11-taskmap}
\end{equation}
For a tool-position task, $f$ is the position part of the direct kinematics.
For a full-pose task it includes orientation as well. Configuration space
and task space are related by this map; neither is the mathematical dual
or set complement of the other.

The \emph{reachable task set} is
\begin{equation}
 \mathcal X_{\mathrm{reach}}=f(\mathcal C).
 \label{eq:l11-reachable}
\end{equation}
It can be smaller than the space in which task coordinates are expressed.
For a position task, this set is the tool's position workspace. A desired
task value may correspond to no configuration, one configuration, or
multiple configurations. The RR examples in Chapter~\ref{chap:lecture10}
already showed that different joint configurations can yield the same
tool position.

\section{An RP arm connects the two descriptions}
\label{sec:l11-rp}
Consider a revolute joint followed by a prismatic joint. The first axis
is perpendicular to the drawing plane; the second lies in that plane and
extends toward the tool. The two axes intersect at the base. Let $q_1$
be the revolute coordinate and $q_2$ the tool's distance along the
prismatic axis, with no additional extension offset.

\begin{figure}[htbp]
\centering\begin{tikzpicture}[scale=1.05]
\coordinate (O) at (0,0); \coordinate (P) at (30:4);
\begin{scope}[rotate=30]
\draw[fill=gray!25,thick] (.25,-.32) rectangle (2.5,.32);
\draw[fill=gray!10,thick] (1.6,-.16) rectangle (4,.16);
\draw[thick] (4,-.45)--(4.35,-.45) (4,-.45)--(4,.45)--(4.35,.45);
\draw[<->] (.4,-.7)--(4,-.7) node[midway,below] {$q_2$};
\end{scope}
\draw[fill=gray!15,thick] (0,0) circle (.23);
\fill (0,0) circle (.05);
\node[below left] at (0,-.25) {$O_0=O_1$};
\node[below left] at (-.15,.15) {$J_1=z_0\ (\odot)$};
\draw[axis] (O)--(2,0) node[right] {$x_0$};
\draw[axis] (O)--(0,2.5) node[above] {$y_0$};
\draw[axis,blue!75!black] (O)--(120:2) node[above] {$x_1$};
\draw[blue!65!black,dashed] (120:2.2)--(300:.8);
\draw[axis,red!75!black] (O)--(30:2.9) node[above left] {$z_1,\ J_2$};
\draw[->] (1,0) arc (0:120:1);
\node at (60:1.23) {$q_1$};
\fill (P) circle (.05);
\draw[axis,blue!75!black] (P)--++(120:1.3) node[above] {$x_2$};
\draw[axis,red!75!black] (P)--++(30:1.5) node[right] {$z_2$};
\node[below right] at (P) {$O_2$};
\node[align=center] at (3.8,-1.15) {$y_1,y_2$ point out of the page\\$x_1$: common normal};
\end{tikzpicture}

\caption{The planar RP arm and its DH frames. The revolute axis $J_1=z_0$
points out of the page; $J_2$ follows the slider. Frames 0 and 1 share an
origin. The common normal is along $x_1$, perpendicular to both joint
axes. This choice makes the tool bearing $q_1-\pi/2$, rather than $q_1$.}
\label{fig:l11-rp}
\end{figure}

\subsection{Frame assignment and direct kinematics}
Choose frame 0 to coincide with the world frame, with $z_0$ out of the
page. Frame 1 has the same origin, $z_1$ along the positive prismatic
direction, and $x_1$ along the common normal. Choose its sense so that
the twist from $z_0$ to $z_1$ about $x_1$ is $+\pi/2$. At the tool,
choose $x_2$ parallel to $x_1$ and $z_2$ parallel to $z_1$.

The standard DH table is
\begin{center}
\begin{tabular}{cccccc}
\toprule
$i$ & Type & $\theta_i$ & $\alpha_i$ & $a_i$ & $d_i$\\
\midrule
1 & R & $q_1$ & $\pi/2$ & $0$ & $0$\\
2 & P & $0$ & $0$ & $0$ & $q_2$\\
\bottomrule
\end{tabular}
\end{center}
With $c_1=\cos q_1$ and $s_1=\sin q_1$,
\[
 H_1^0=\begin{bmatrix}c_1&0&s_1&0\\s_1&0&-c_1&0\\
 0&1&0&0\\0&0&0&1\end{bmatrix},\qquad
 H_2^1=\begin{bmatrix}1&0&0&0\\0&1&0&0\\0&0&1&q_2\\0&0&0&1\end{bmatrix}.
\]
Their product gives
\begin{equation}
 H_2^0=\begin{bmatrix}
 c_1&0&s_1&q_2s_1\\s_1&0&-c_1&-q_2c_1\\
 0&1&0&0\\0&0&0&1
 \end{bmatrix}.
 \label{eq:l11-rp-H}
\end{equation}
The translation column is the tool position; the upper-left block is its
orientation. For the planar position task, retain only
\begin{equation}
 \boxed{\xi=\begin{bmatrix}x\\y\end{bmatrix}
 =f(q)=\begin{bmatrix}q_2\sin q_1\\-q_2\cos q_1\end{bmatrix}.}
 \label{eq:l11-rp-map}
\end{equation}
The sine and negative cosine follow from the chosen DH frames. At $q_1=0$,
the slider points along $-y_0$; at $q_1=\pi/2$, it points along $+x_0$.
Thus $q_1$ is the angle of $x_1$ from $x_0$, not the polar bearing of
the extension itself.

\subsection{Joint limits and boundary points}
Assume the limits
\begin{equation}
 0\le q_1\le\frac{3\pi}{2},\qquad
 1\,\mathrm m\le q_2\le2\,\mathrm m.
 \label{eq:l11-limits}
\end{equation}
The configuration space is a rectangle in the $(q_1,q_2)$ coordinate
plot. Its horizontal axis is measured in radians and its vertical axis
in meters; these are not spatial $x$ and $y$ axes.

Mapping selected boundary points produces the following correspondence:
\begin{center}
\begin{tabular}{ccccc}
\toprule
Point & $q_1$ (rad) & $q_2$ (m) & $x$ (m) & $y$ (m)\\
\midrule
A & $0$ & $1$ & $0$ & $-1$\\
B & $0$ & $2$ & $0$ & $-2$\\
C & $3\pi/2$ & $1$ & $-1$ & $0$\\
D & $3\pi/2$ & $2$ & $-2$ & $0$\\
E & $\pi$ & $1$ & $0$ & $1$\\
F & $\pi$ & $2$ & $0$ & $2$\\
L & $\pi/2$ & $1$ & $1$ & $0$\\
M & $\pi/2$ & $2$ & $2$ & $0$\\
\bottomrule
\end{tabular}
\end{center}
These points identify useful corners and directions, but the continuous
boundary follows from the map itself. Since
$x^2+y^2=q_2^2$, constant extension traces a circular arc. Constant $q_1$
traces a radial segment. With polar bearing $\phi=q_1-\pi/2$, the reachable
set is
\begin{equation}
 \mathcal X_{\mathrm{reach}}=
 \{(r\cos\phi,r\sin\phi):1\le r\le2,\ -\pi/2\le\phi\le\pi\},
 \label{eq:l11-sector}
\end{equation}
where lengths are in meters. This is three quarters of an annulus: the
lower-left sector is absent because of the revolute joint limit.

\begin{figure}[htbp]
\centering\begin{tikzpicture}[scale=1.05]
\filldraw[fill=blue!8,draw=blue!55] (0,-2) arc(-90:180:2)--(-1,0) arc(180:-90:1)--cycle;
\draw[axis] (-2.5,0)--(2.6,0) node[right] {$x$};
\draw[axis] (0,-2.5)--(0,2.6) node[above] {$y$};
\foreach \x/\y/\l/\a in {0/-1/A/right,0/-2/B/right,-1/0/C/below,-2/0/D/below,0/1/E/left,0/2/F/left,1/0/L/below,2/0/M/below}
 {\fill (\x,\y) circle(.045);\node[\a] at (\x,\y) {\l};}
\node at (0,-3) {Reachable position (m)};
\begin{scope}[shift={(5,0)},x=1cm,y=1.5cm]
\filldraw[fill=blue!8,draw=blue!55] (0,1) rectangle (3,2);
\draw[axis] (-.3,0)--(3.7,0) node[right] {$q_1$};
\draw[axis] (0,-.2)--(0,2.5) node[above] {$q_2$ (m)};
\foreach \x/\t in {1/{\pi/2},2/{\pi},3/{3\pi/2}} {
\draw[projection] (\x,0)--(\x,2);\node[below] at (\x,0) {$\t$};}
\foreach \y in {1,2} \node[left] at (0,\y) {$\y$};
\foreach \x/\y/\l/\a in {0/1/A/below right,0/2/B/above right,3/1/C/below,3/2/D/above,2/1/E/below,2/2/F/above,1/1/L/below,1/2/M/above}
 {\fill (\x,\y) circle(.04);\node[\a] at (\x,\y) {\l};}
\node at (1.5,-2) {Configuration coordinates};
\end{scope}
\end{tikzpicture}

\caption{The same RP model represented in task coordinates (left) and
configuration coordinates (right). Matching letters denote corresponding
points. The nonlinear map converts horizontal edges of the rectangle into
circular arcs and vertical edges into radial segments.}
\label{fig:l11-mapped}
\end{figure}

\section{A straight line depends on the chosen space}
\label{sec:l11-straight}
Take a starting configuration $q_S=(\pi/2,1.5)$ and a goal configuration
$q_G=(3\pi/2,1.5)$, with the extension in meters. Their task positions
are $\xi_S=(1.5,0)$ and $\xi_G=(-1.5,0)$.

A straight segment between these task points passes through the origin,
where $q_2=0$ would be needed. It therefore leaves the reachable set.
Its endpoints are reachable, but its intermediate points are not all
reachable.

Instead, interpolate the configuration coordinates:
\[
 q(\lambda)=\begin{bmatrix}\pi/2+\pi\lambda\\1.5\end{bmatrix},
 \qquad0\le\lambda\le1.
\]
Every point of this horizontal segment stays within the joint limits.
The resulting tool path is
\begin{equation}
 f(q(\lambda))=
 \begin{bmatrix}1.5\cos(\pi\lambda)\\1.5\sin(\pi\lambda)\end{bmatrix},
 \label{eq:l11-arc}
\end{equation}
the upper semicircle of radius $1.5\,\mathrm m$.

\begin{figure}[htbp]
\centering\begin{tikzpicture}[scale=1.05]
\filldraw[fill=blue!8,draw=blue!40] (0,-2) arc(-90:180:2)--(-1,0) arc(180:-90:1)--cycle;
\draw[axis] (-2.4,0)--(2.5,0) node[right] {$x$};
\draw[axis] (0,-2.3)--(0,2.5) node[above] {$y$};
\draw[red!75!black,very thick,dashed,->] (1.5,0)--(-1.5,0);
\draw[green!45!black,very thick,->] (1.5,0) arc(0:180:1.5);
\fill (1.5,0) circle(.05) node[below] {$S$};
\fill (-1.5,0) circle(.05) node[below] {$G$};
\node at (0,-2.8) {Task space};
\begin{scope}[shift={(5,0)},y=1.5cm]
\filldraw[fill=blue!8,draw=blue!40] (0,1) rectangle (3,2);
\draw[axis] (0,0)--(3.7,0) node[right] {$q_1$};
\draw[axis] (0,0)--(0,2.5) node[above] {$q_2$};
\foreach \x/\t in {1/{\pi/2},2/{\pi},3/{3\pi/2}} {
\draw[projection] (\x,0)--(\x,2);\node[below] at (\x,0) {$\t$};}
\foreach \y in {1,1.5,2} \node[left] at (0,\y) {$\y$};
\draw[very thick,green!45!black,->] (1,1.5)--(3,1.5);
\fill (1,1.5) circle(.05) node[above] {$S$};
\fill (3,1.5) circle(.05) node[above left] {$G$};
\node at (1.5,-1.87) {Configuration space};
\end{scope}
\end{tikzpicture}

\caption{The dashed task-space segment crosses unreachable positions.
The solid configuration-space segment maps to the solid circular tool
path. Straight interpolation in one space need not be straight in the
other. This comparison assumes no collision obstacles.}
\label{fig:l11-paths}
\end{figure}

The rectangle of independent joint intervals is convex: every segment
between its points remains in it. This property applies to this coordinate
domain, not to every configuration space or free space. Periodic joints,
closed-chain constraints, and obstacles require additional care. Likewise,
being inside the joint limits does not establish collision freedom or
dynamic feasibility.

\section{Obstacles exclude configurations}
An obstacle occupies a region $\mathcal O$ of physical space. Let
$\mathcal A(q)$ be the physical region occupied by the robot at configuration
$q$. With contact treated as collision, the configuration-space obstacle is
\begin{equation}
 \mathcal C_{\mathrm{obs}}=
 \{q\in\mathcal C:\mathcal A(q)\cap\mathcal O\ne\varnothing\},
 \qquad
 \mathcal C_{\mathrm{free}}=\mathcal C\setminus\mathcal C_{\mathrm{obs}}.
 \label{eq:l11-cobs}
\end{equation}
Self-collision constraints can be excluded in the same way. In this
representation the robot configuration is a point, while the robot's
shape is incorporated into the forbidden set.

For an idealized point tool, one can instead test $f(q)\in\mathcal O$.
For example, a disk of radius $\rho$ centered at $(x_o,y_o)$ excludes
the RP tool configurations satisfying
\[
 (q_2\sin q_1-x_o)^2+(-q_2\cos q_1-y_o)^2\le\rho^2.
\]
This generally produces a curved boundary in $(q_1,q_2)$ coordinates.
It is only a tool-point collision test: a link or slider can collide
even when its endpoint lies outside the obstacle. Full-arm planning must
use the robot's occupied geometry, not just this task map.

Planning in configuration space checks candidate paths against
$\mathcal C_{\mathrm{free}}$. Planning in task space must additionally
ensure a consistent, admissible configuration exists along the task path.
The choice of representation can simplify a problem, but does not remove
its reachability, collision, or motion constraints.

\clearpage
\section{Additional exercises}
\begin{enumerate}
\item \textbf{Configuration or task?} Give a configuration vector and
a position-only task vector for a planar mobile robot and for a fixed-base
RR arm. Explain why their dimensions can differ and why configuration
alone is not a complete dynamic state.

\item \textbf{Selecting outputs.} Specify suitable task coordinates for
altitude hold, positioning a gripper, and keeping a container upright.
Which orientation components can be left unspecified in each case?

\item \textbf{RP direct kinematics.} Reconstruct the two local DH matrices
and multiply them to obtain Equation~\eqref{eq:l11-rp-H}. Identify the
directions of $x_2$ and $z_2$ at $q_1=0$ and $q_1=\pi/2$.

\item \textbf{Mapping a configuration.} For $q_1=\pi/4$ and
$q_2=1.5\,\mathrm m$, calculate the RP tool position. State its distance
and polar bearing from the base. Explain the $\pi/2$ offset between the
joint coordinate and the tool bearing.

\item \textbf{Changed limits.} Replace the RP limits by
$0\le q_1\le\pi$ and $0.5\le q_2\le1.5$ meters. Sketch the configuration
rectangle and its reachable task set. Map all four corners and identify
the inner and outer arcs.

\item \textbf{Comparing paths.} For the start and goal in
Section~\ref{sec:l11-straight}, calculate both path midpoints: the midpoint
of the task-space segment and the image of the configuration-space
midpoint. Which is reachable under the stated limits? Give a reason why
neither calculation alone establishes full-arm collision freedom.

\item \textbf{A point-tool obstacle.} For the RP model, take a disk
centered at $(1.8,0)\,\mathrm m$ with radius $0.2\,\mathrm m$. At
$q_1=\pi/2$, find the forbidden interval of $q_2$ under the point-tool
model, treating contact as collision. Explain why a full-link collision
test could forbid additional configurations.
\end{enumerate}
